% Préambule commun aux manuels de Tutti (français et anglais).
% Compilation : pdflatex manuel-fr.tex (deux fois), idem manuel-en.tex, depuis docs/src/.
\documentclass[11pt,a4paper]{article}
\usepackage[utf8]{inputenc}
\usepackage[T1]{fontenc}
\usepackage{lmodern}
\usepackage{microtype}
\usepackage[a4paper,margin=2cm,headheight=24pt]{geometry}
\usepackage{graphicx}
\usepackage[table]{xcolor}
\usepackage[most]{tcolorbox}
\usepackage{enumitem}
\usepackage{booktabs}
\usepackage{array}
\usepackage{tabularx}
\usepackage[nobottomtitles*]{titlesec}
\usepackage{needspace}
\usepackage{fancyhdr}
\usepackage{fontawesome5}
\usepackage{etoolbox}
\DeclareUnicodeCharacter{2192}{\ensuremath{\rightarrow}}
\DeclareUnicodeCharacter{2191}{\ensuremath{\uparrow}}
\DeclareUnicodeCharacter{2193}{\ensuremath{\downarrow}}
\DeclareUnicodeCharacter{25B2}{\ensuremath{\blacktriangle}}
\DeclareUnicodeCharacter{25BC}{\ensuremath{\blacktriangledown}}
\DeclareUnicodeCharacter{21BA}{\ensuremath{\circlearrowleft}}
\DeclareUnicodeCharacter{21BB}{\ensuremath{\circlearrowright}}
\usepackage{amssymb}

% Palette de l'application : ivoire, bordeaux, or
\definecolor{bordeaux}{HTML}{2C6E49}
\definecolor{bordeauxl}{HTML}{E4F0E8}
\definecolor{or}{HTML}{B8860B}
\definecolor{orl}{HTML}{FBF3DC}
\definecolor{sapin}{HTML}{2C7A55}
\definecolor{sapinl}{HTML}{E2F1E9}
\definecolor{sky}{HTML}{1C6DB4}
\definecolor{skyl}{HTML}{E6F0F8}
\definecolor{summit}{HTML}{6B7A8F}
\definecolor{night}{HTML}{2A2420}
\definecolor{ink}{HTML}{1F2430}
\AtBeginDocument{\color{ink}}

\titleformat{\section}
  {\Large\bfseries\sffamily\color{bordeaux}}
  {\makebox[1.6em][l]{\thesection}}{0.3em}{}
  [{\vspace{1pt}\color{bordeaux!35}\titlerule[1.3pt]}]
\titlespacing*{\section}{0pt}{14pt}{7pt}
\titleformat{\subsection}{\large\bfseries\sffamily\color{night}}{\thesubsection}{0.5em}{}
\pretocmd{\section}{\needspace{10\baselineskip}}{}{}
\pretocmd{\subsection}{\needspace{7\baselineskip}}{}{}
\widowpenalty=10000
\clubpenalty=10000
\raggedbottom
\setlength{\parindent}{0pt}
\setlength{\parskip}{6pt}
\setlist[enumerate]{itemsep=3pt,leftmargin=6mm}
\setlist[itemize]{itemsep=2pt,leftmargin=5mm}

\pagestyle{fancy}
\fancyhf{}
\renewcommand{\headrulewidth}{0.4pt}
\renewcommand{\headrule}{\hbox to\headwidth{\color{bordeaux!40}\leaders\hrule height \headrulewidth\hfill}}
\fancyfoot[C]{\sffamily\small\color{summit}\thepage}
\renewcommand{\sectionmark}[1]{\markboth{#1}{}}

\newcommand{\menu}[1]{\textsf{\textbf{\textcolor{bordeaux}{#1}}}}
\newcommand{\touche}[1]{\tcbox[on line,boxsep=0pt,left=2pt,right=2pt,top=1.5pt,bottom=1.5pt,
  colback=white,colframe=summit,boxrule=0.5pt,arc=1.5pt]{\footnotesize\textsf{#1}}}
\newcommand{\fichier}[1]{\texttt{\small #1}}

\newcommand{\boite}[5]{\newtcolorbox{#1}{enhanced,unbreakable,colback=#2,colframe=#3,
  coltitle=white,fonttitle=\bfseries\sffamily\small,boxrule=0.9pt,arc=2.5pt,
  left=2.5mm,right=2.5mm,top=1.2mm,bottom=1.2mm,title={#4~~#5}}}
\newcommand{\capture}[2][1]{\par\needspace{8\baselineskip}\begin{center}\setlength{\fboxsep}{0pt}\setlength{\fboxrule}{0.6pt}%
  \fcolorbox{summit!50}{white}{\includegraphics[width=\dimexpr#1\linewidth-1.3pt\relax]{#2}}\end{center}}

\usepackage[french]{babel}
\usepackage[colorlinks=true,linkcolor=bordeaux,urlcolor=bordeaux]{hyperref}
\hypersetup{pdftitle={Da Capo — manuel d'utilisation},pdfauthor={},pdfcreator={Da Capo}}
\boite{astucebox}{orl}{or!85!black}{\faLightbulb}{Bon à savoir}
\boite{securitebox}{sapinl}{sapin}{\faShield*}{Vos partitions sont en sécurité}
\boite{attentionbox}{alertl}{alert}{\faExclamationTriangle}{Attention}
\newcommand{\astuce}[1]{\par\needspace{5\baselineskip}\begin{astucebox}#1\end{astucebox}}
\newcommand{\attention}[1]{\par\needspace{5\baselineskip}\begin{attentionbox}#1\end{attentionbox}}
\definecolor{alert}{HTML}{B5382B}
\definecolor{alertl}{HTML}{F8E6E3}
\fancyhead[L]{\sffamily\small\color{bordeaux}\textbf{DA CAPO — MANUEL}}
\fancyhead[R]{\sffamily\small\color{summit}\leftmark}
\input{version}

\begin{document}

\begin{titlepage}
\centering
\vspace*{2.2cm}
\includegraphics[width=3.6cm]{../../logo_1024.png}\par
\vspace{8mm}
{\sffamily\Huge\bfseries\color{bordeaux} DA CAPO\par}
\vspace{3mm}
{\sffamily\LARGE\color{ink} Sortir ses partitions de Finale\par}
\vspace{2mm}
{\sffamily\large\color{summit} Manuel d'utilisation \textbullet{} version \versiondacapo\par}
\vspace{12mm}
\begin{tcolorbox}[enhanced,width=0.82\linewidth,colback=orl,colframe=or,arc=4pt,boxrule=1.2pt,left=5mm,right=5mm,top=3mm,bottom=3mm]
\sffamily\centering
\textit{Da capo} : « depuis le début », en musique.\\[2mm]
\raggedright
\textbf{1.} Dossiers de partitions \quad \textbf{2.} Dossier de travail \quad \textbf{3.} Finale \quad \textbf{4.} Commencer
\end{tcolorbox}
\vfill
{\footnotesize\color{summit}\sffamily Windows \textbullet{} macOS \textbullet{} Linux\par}
\end{titlepage}

\tableofcontents
\newpage

\section{Pourquoi Da Capo}

Finale n'est plus développé depuis 2024. Il fonctionne encore, mais plus aucune mise à jour ne
viendra le réparer si un jour Windows le casse. Or ses fichiers \fichier{.mus} et \fichier{.musx}
\textbf{ne s'ouvrent que dans Finale}. La porte de sortie s'appelle \textbf{MusicXML} : un format
ouvert, que lisent MuseScore Studio (libre et gratuit), Dorico et Sibelius.

Finale sait exporter une partition en MusicXML, mais une à la fois. Avec des centaines de
partitions, c'est des jours de clics. \textbf{Da Capo fait ces clics à votre place}, partition après
partition, pendant que vous faites autre chose — et vérifie le résultat.

\begin{center}
\renewcommand{\arraystretch}{1.35}
\begin{tabularx}{\linewidth}{>{\bfseries\color{bordeaux}}l X}
\toprule
1. Trouver & toutes les partitions Finale des dossiers choisis, sous-dossiers compris. \\
2. Copier & chacune dans le dossier de travail, copie vérifiée octet par octet. \\
3. Exporter & Finale ouvre chaque \emph{copie}, l'exporte en MusicXML (et en PDF de référence), la referme. \\
4. Vérifier & chaque MusicXML est relu : parties, mesures, notes, paroles, nuances. \\
5. Préparer & chaque MusicXML devient une partition MuseScore (\fichier{.mscz}) + un PDF pour comparer. \\
6. Rendre compte & \fichier{rapport.html} : ce qui a marché, ce qui a échoué, ce qui mérite un coup d'œil. \\
\bottomrule
\end{tabularx}
\end{center}

\begin{securitebox}
\begin{itemize}
\item Les fichiers d'origine ne sont \textbf{jamais ouverts ni modifiés} : Finale ne voit que des copies.
\item Les copies sont vérifiées (empreinte SHA-256) et mises en lecture seule.
\item Rien n'est envoyé sur Internet : tout se passe sur votre ordinateur.
\end{itemize}
\end{securitebox}

\section{Avant de commencer}

\begin{enumerate}
\item \textbf{Installez Da Capo} : \fichier{DaCapo-Setup.exe} sous Windows (aucun mot de passe
      administrateur). Au premier lancement, Windows peut afficher « Windows a protégé votre
      ordinateur » : \menu{Informations complémentaires} → \menu{Exécuter quand même}.
\item \textbf{Installez MuseScore Studio} (facultatif mais conseillé) depuis \fichier{musescore.org}
      — et non \fichier{musescore.com}, qui est un site de partitions en partie payant. Le logiciel
      est gratuit et le reste.
\item \textbf{Ouvrez une fois à la main le Finale que Da Capo utilisera}, surtout s'il ne sert jamais :
      réglez ce qu'il demande au premier lancement (activation, mise à jour, fenêtre d'accueil), puis
      \textbf{refermez-le}.
\item \textbf{Avant Finale 25}, le PDF de référence passe par l'imprimante virtuelle
      \textbf{CutePDF Writer} (gratuite) : vérifiez qu'elle est installée.
\end{enumerate}

\astuce{Si plusieurs Finale sont installés, Da Capo propose le plus récent : il ouvre sans
difficulté les fichiers des anciennes versions et écrit un MusicXML bien plus fidèle (version 4.0
pour Finale 27, contre un « MusicXML de base » pour Finale 2003 à 2010).}

\section{Les quatre étapes}

\capture[0.78]{../capture_prete-fr.png}

\begin{description}[style=nextline,leftmargin=5mm,font=\sffamily\bfseries\color{bordeaux}]
\item[1. Où sont vos partitions Finale ?]
\menu{Ajouter un dossier…}, ou glissez le dossier depuis l'Explorateur. Tous les sous-dossiers sont
parcourus. Da Capo compte aussitôt les partitions trouvées (\fichier{.mus} jusqu'à Finale 2011,
\fichier{.musx} après). Les copies de secours de Finale (« Sauvegarde de… », « Backup of… ») sont
ignorées.
\item[2. Où ranger le résultat ?]
Un dossier vide, sur un disque qui a de la place (prévoyez trois à quatre fois la taille des
partitions). Tout y sera rangé, dans la même arborescence que vos dossiers d'origine.
\item[3. Avec quel Finale ?]
Le plus récent est proposé et recommandé. La première fois, cliquez sur \menu{Diagnostic…}
(section~\ref{sec:diagnostic}).
\item[4. Que produire ?]
Le \textbf{MusicXML}, toujours. Le \textbf{PDF de référence} : la partition exactement comme Finale
l'imprime — c'est votre étalon pour vérifier, et une archive qui s'ouvrira toujours. Les
\textbf{fichiers MuseScore}, si MuseScore Studio est installé.
\end{description}

Cliquez ensuite sur \menu{Commencer}.

\section{Pendant l'export}

\capture[0.78]{../capture_export-fr.png}

Finale s'ouvre et travaille seul : on le voit ouvrir chaque partition, exporter, refermer.
\textbf{Ne touchez ni à la souris ni au clavier} : un clic au mauvais moment détournerait une boîte
de dialogue. Comptez de quelques secondes à une minute par partition ; pour plusieurs centaines,
lancez le travail le soir et laissez tourner la nuit. L'ordinateur ne se met pas en veille.

\begin{itemize}
\item \menu{Arrêter} termine la partition en cours puis s'arrête proprement.
\item Le bouton devient alors \menu{Reprendre} : Da Capo saute tout ce qui est déjà fait.
      Même après une coupure de courant, il ne refait que la partition en cours.
\item Si Finale affiche un message (fichier ancien à convertir, police manquante…), Da Capo le lit,
      le note dans le rapport, puis le referme. À la question « Enregistrer les modifications ? »,
      il répond toujours \textbf{Non}.
\item Si Finale ne répond plus, Da Capo l'arrête, le relance et passe à la partition suivante ;
      la partition rebelle est marquée en échec et sera retentée à la reprise.
\end{itemize}

\section{Le dossier de travail}

\begin{center}
\renewcommand{\arraystretch}{1.3}
\begin{tabularx}{\linewidth}{>{\ttfamily\small}l X}
\toprule
originaux/ & les copies de vos partitions, en lecture seule \\
musicxml/ & les MusicXML exportés par Finale — à garder précieusement \\
pdf-finale/ & les PDF de référence \\
musescore/ & les partitions MuseScore (\fichier{.mscz}) et leur PDF \\
rapport.html & le bilan, à ouvrir dans le navigateur \\
\_dacapo/ & le journal qui permet de reprendre, et le relevé du diagnostic \\
\bottomrule
\end{tabularx}
\end{center}

Le dossier entier peut être déplacé ou copié (sur un disque de sauvegarde, par exemple) : les liens
du rapport suivent.

\section{Le rapport}

\menu{Ouvrir le rapport} affiche une page avec, en haut, le bilan (partitions, MusicXML, PDF,
MuseScore, échecs), puis une ligne par partition :

\begin{center}
\renewcommand{\arraystretch}{1.35}
\begin{tabularx}{\linewidth}{l X}
\toprule
\textcolor{alert}{\rule{4pt}{9pt}}\ \textbf{rouge} & l'export a échoué ; la raison est écrite (et les messages de Finale) \\
\textcolor{or}{\rule{4pt}{9pt}}\ \textbf{orange} & exporté, mais à vérifier : partition vide, message de Finale pendant l'export \\
\textcolor{sapin}{\rule{4pt}{9pt}}\ \textbf{vert} & tout s'est bien passé \\
\bottomrule
\end{tabularx}
\end{center}

Pour chaque partition : le nombre de parties, de mesures et de notes trouvées dans le MusicXML, et
des liens vers chaque fichier produit.

\section{Vérifier une partition}

Ouvrez côte à côte le \textbf{PDF de référence} (Finale) et le \textbf{PDF MuseScore}. Les notes,
les rythmes, les paroles, les nuances et les parties passent bien en MusicXML. Ce qui demande
souvent une retouche dans MuseScore :
\begin{itemize}
\item la \textbf{mise en page} : sauts de système et de page, espacements, taille de la portée ;
\item les \textbf{expressions personnalisées} et symboles rares, les polices musicales particulières ;
\item la \textbf{forme des liaisons} et la position fine des textes ;
\item ce qui ne concerne que la lecture audio de Finale (Human Playback).
\end{itemize}
Commencez par les partitions dont vous vous servez encore ; les autres restent archivées en trois
exemplaires (original, MusicXML, PDF).

\attention{Gardez les originaux \fichier{.mus} et les copies du dossier de travail. Un jour,
un meilleur convertisseur pourra refaire l'export à partir d'eux.}

\section{Le diagnostic}\label{sec:diagnostic}

Chaque version de Finale range ses commandes un peu différemment. Da Capo les retrouve par leur nom,
en français comme en anglais. Le bouton \menu{Diagnostic…} le vérifie sur votre ordinateur : il
ouvre Finale avec la copie d'une partition, relève ses menus et ses fenêtres, puis referme tout,
sans rien modifier.

\capture[0.78]{../capture_diagnostic-fr.png}

Si tout est reconnu, c'est prêt. Sinon, \menu{Envoyer le relevé} le transmet au développeur, qui
adaptera Da Capo à votre Finale ; le relevé est aussi enregistré dans \fichier{\_dacapo}.

\section{Problèmes fréquents}

\begin{center}
\renewcommand{\arraystretch}{1.35}
\begin{tabularx}{\linewidth}{>{\raggedright\arraybackslash}p{0.38\linewidth} X}
\toprule
\textbf{Message} & \textbf{Que faire} \\
\midrule
Finale ne démarre pas & L'ouvrir une fois à la main (activation, mise à jour), le refermer, recommencer. \\
Commande introuvable dans les menus & Faire un diagnostic et l'envoyer. \\
Imprimante CutePDF Writer introuvable & L'installer (gratuite), ou décocher le PDF de référence. \\
Finale n'a pas ouvert la partition à temps & Fichier abîmé ou très gros : l'ouvrir à la main dans Finale pour voir. \\
MuseScore a mis trop longtemps & Ouvrir le MusicXML directement dans MuseScore. \\
\bottomrule
\end{tabularx}
\end{center}

\section{Mises à jour et confidentialité}

Au démarrage, Da Capo vérifie discrètement s'il existe une nouvelle version et propose de
l'installer. Avec votre accord, un rapport est envoyé s'il rencontre un problème : version du
programme et du système, erreur rencontrée ; les noms de vos fichiers et votre nom d'utilisateur
sont masqués. Aucune partition n'est jamais envoyée.

Sous macOS et Linux, Finale ne peut pas être piloté : Da Capo y copie les partitions, vérifie des
MusicXML exportés ailleurs et prépare les fichiers MuseScore.

\vfill
{\footnotesize\color{summit}\sffamily
Finale est une marque de MakeMusic, MuseScore de Muse Group ; Da Capo n'est lié à aucun des deux.
Polices Source Sans 3 et EB Garamond (SIL Open Font License). Exemples de ce manuel : fichiers fictifs.\par}

\end{document}
